\section{Discussion and directions for further tests}
In both experiments, an informative control uses two histories in which the same observation calls for different answers. A system that sees only what the histories share cannot answer correctly in both. Such a pair establishes what information is necessary. To attribute use of that information to a particular system, one must also show what happens when its carrier is delivered, replaced or removed while the subsequent input is held fixed. Geometric coverage of a cone, a correlate recoverable from the data, and high mean accuracy on their own do not establish such an effect.

\paragraph{Scope of the conclusions.}
The relationship between zero-error coding and graph coloring is known in information theory \cite{witsenhausen,charpenay}. The contribution here lies in specifying testable interventions, providing concrete discrete witnesses and restricting the channel between two episodes. On the infinite line, for every time $t$ and each coordinate $j$ inside the window, there exists a pair of configurations that differ only at $j$ and give different target bits. Outside the window, changing $j$ alone preserves the target bit for every initial configuration. In the OLMo experiments, each table row concerns its own combination of frozen checkpoint (Base or Instruct) and interface on one finite panel. Whether the system has subjective experience or lacks it is beyond the scope of these data.

\paragraph{Transfer to other domains.}
In a mechanistic audit, first establish that the model can solve the task with the history intact. One can then transfer a candidate hidden carrier and test what happens when it is erased. In consciousness research, including the discussion of disorders of consciousness, it is useful to distinguish availability of a report, assumptions about experience and a decision to take precautions. Clinical application of that distinction requires independent validation. The question of a state's provenance also carries over to digital forensics and distributed computing. A correct recording can reproduce a live response to a fixed query. In a separate series with a change of address, the live channel scored 8/8 and the prerecorded tape 6/8. The origin of a branch and its ability to answer a new query can be tested independently.

\paragraph{Addressable feedback in OLMo (T3, \texttt{NOT\_RUN}).}
Once testing on panel 33--44 is complete, this experiment requires a separate panel. The model retains two of its earlier choices $h_0,h_1$. From the content of a new task it must infer which choice needs correction. The reference address $K$ is hidden. Before seeing the string $\mathbf{x}$, the model produces its own query bit $q$. The observer holds two feedback bits $e_0,e_1$ and returns one of them. The model must then apply the target rule $h_K\oplus e_K$ to the string. The observer's reply is the model's only source of information about the feedback bits.

For the same histories, deliver to the observer address $q$ (live query), $0$ (fixed address) or $1-q$ (swapped address). In every condition the recipient sees one bit at the same position. Log the delivered address and the swap indicator and keep both hidden from the recipient. Measure selection of the useful query ($q=K$) separately from final-answer accuracy when $L\ne R$. A paired comparison at $e_0\ne e_1$ shows whether the answer changed when the address was swapped.

On strings with differing endpoints, the two values of $K$ occur equally often. For each task, pair of past choices and string, all four feedback pairs occur equally often. The accuracy ceiling with a fixed address is $75\%$. When $K=0$, the reply $e_0$ gives the required feedback. When $K=1$, the two values of $e_1$ remain equiprobable even after receipt of $e_0$. Optimal accuracy is therefore $\tfrac12\cdot1+\tfrac12\cdot\tfrac12=\tfrac34$. Whether a live query exceeds this bound is tested on the entire balanced panel of strings with differing endpoints.

\begin{center}\small
\begin{tabularx}{\linewidth}{@{}p{.23\linewidth}X@{}}\toprule
Task at this snapshot & Evidence required \\\midrule
G4-1, panel 33--44 & Freeze the two interfaces being compared on dev before first opening their shared panel. Exact one-bit transfer requires all 12/12 complete blocks, together with the erasure, format and channel-boundary controls.\\
G3, panel 21--32 & After qualifying the textual channel, replace the text with one prespecified activation pulse. Turning the pulse off must make the two feedback branches identical for the same choice and input.\\
AI.04: multidimensional G4 and DAS & Transfer a specified internal state with the source positions inaccessible. To localize a subspace, prespecify the layer, mediation metric, tolerance and an independent control. $C_{\rm med}$ also requires another task using the same rule.\\
R110.11, slot A6 & Implement storage and return entirely in homogeneous binary Rule~110. The demonstrated autonomous construction uses an extended alphabet.\\
Formalizing merger and return & Formalize the exhaustive check of 4096 environments and the constructive witness for $t=4\to5$ in Lean. Verification of the root geometry theorem and 11/11 queries in the observer module has been reproduced (Appendix~\ref{app:lean}).\\
$D_{\rm dyn}$ & Show a change in organization that persists beyond the episode, and its effect on a new input under controlled reset.\\
\bottomrule\end{tabularx}
\end{center}
The gate protocol in Appendix~\ref{app:gates} specifies the outcomes under which the proposed channel fails its test. A negative result for a particular channel limits what can be inferred about the chosen carrier. Appendix~\ref{app:registers} lists the associated hypotheses and their statuses.
